\documentclass{situationist-preprint}

\title{Population Ethics Without Paradox}
\msid{ETH-09-c}
\msfamily{ETH-09}
\msdiscipline{Normative and Applied Ethics}
\mstype{Consequence}
\mskeywords{population ethics; mere addition paradox; sadistic conclusion; procreative asymmetry; impossibility theorems}

\begin{document}
\maketitle

\begin{abstract}
Population ethics is often presented as a field of paradoxes: sets of individually plausible principles that cannot all be true. ETH-09-a argues that the repugnant conclusion is true and not repugnant. We show what follows for the other standard puzzles. The mere addition paradox dissolves, since its conclusion is the repugnant conclusion. The impossibility results lose their force, since the condition the total view violates in them is the demand to avoid that conclusion. The sadistic conclusion is avoided outright. The procreative asymmetry is explained, in the part that is true, by the distinction between impersonal reasons and obligations, and the non-identity problem is handled by ETH-08. What remains is a field without paradox, though not without hard cases.
\end{abstract}

\section{What a paradox is}

By a \emph{paradox} in population ethics we mean a set of principles, each of which has strong intuitive support, that are jointly inconsistent. A paradox is resolved when one of its principles is shown to lack the support it seemed to have. A theory that resolves a paradox in this way may still have consequences that are surprising; surprise is not paradox.

ETH-09-a establishes the repugnant conclusion from transitivity, mere addition and non-anti-egalitarianism (Assumptions~2.4--2.6 and Theorem~4.2), shows that the total view satisfies these assumptions (Proposition~4.3), and argues that the intuition against the conclusion has no evidential weight (Proposition~6.3). We take this as given and work, where we need a specific axiology, with the total view of Definition~2.3 of ETH-09-a. ETH-09-a does not establish the total view, and the results below that depend on it are marked.

\section{The puzzles}

\paragraph{The mere addition paradox.} \citet[ch.~19]{Parfit1984} presented three judgements that seem jointly inconsistent with the denial of the repugnant conclusion: that A+ is not worse than A, that B$-$ is better than A+, and that Z is not better than A. The first two are instances of Assumptions~2.5 and~2.6 of ETH-09-a, and with transitivity they entail that a sequence of such steps leads to Z. The third is the denial of the repugnant conclusion. Since ETH-09-a shows that the third lacks evidential weight, the paradox is resolved by rejecting it, and the steps of the sequence stand.

\paragraph{Impossibility theorems.} \citet{Arrhenius2000} proved that no axiology satisfies a set of adequacy conditions each of which has seemed compelling. Results of this kind convert the mere addition paradox into a general constraint: every axiology accepts some counterintuitive conclusion. Among the conditions, the one the total view violates is a condition that excludes the repugnant conclusion or a strengthened form of it. If that condition lacks the support it seemed to have, the theorem is not a paradox but a proof that the other conditions, which the total view satisfies, entail the repugnant conclusion. This is the reading that \citet{Zuber2021} recommend to the field when they argue that implying the repugnant conclusion should not be treated as a decisive objection to a population axiology.

\paragraph{The sadistic conclusion.} An axiology implies the sadistic conclusion if, for some population, adding a number of lives with negative value is better than adding a number of lives with positive value. Variable value and critical-level views imply it, as ETH-09-a notes in its Section~5. The total view does not.

\begin{proposition}\label{p:sadistic}
On the total view, for any population $P$, adding any nonempty group of lives at positive levels to $P$ yields a better population than adding any nonempty group of lives at negative levels to $P$.
\end{proposition}

\begin{proof}
Let $P^{+}$ be $P$ with $k > 0$ lives added at positive levels, and $P^{-}$ be $P$ with $j > 0$ lives added at negative levels. Write $T$ for total value. Then $T(P^{+}) = T(P) + s^{+}$ with $s^{+} > 0$, and $T(P^{-}) = T(P) + s^{-}$ with $s^{-} < 0$. So $T(P^{+}) > T(P) > T(P^{-})$, and by Definition~2.3 of ETH-09-a $P^{+}$ is better than $P^{-}$.
\end{proof}

The proposition shows that the sadistic conclusion is not a cost that every view escaping one horn must pay. It is the cost of escaping the repugnant conclusion by discounting additional good lives. A view that accepts the repugnant conclusion has no need to pay it.

\paragraph{The procreative asymmetry.} Many people hold that there is a strong reason not to create a person whose life would be miserable, but no reason to create a person whose life would be happy \citep{Narveson1967}. On the impersonal axiology of ETH-09-a, the second half is false: the addition of a good life makes an outcome better, as Assumption~2.5 of ETH-09-a requires, and so there is an impersonal reason to bring it about. ETH-20 rejects the stronger asymmetry of \citet{Benatar2006} on the same ground.

What is true in the asymmetry is a claim about obligation, not about reasons. ETH-08-a argues that impersonal obligations are not obligations to maximize: the step from value to obligation is governed by deontic principles and by the autonomy of procreative decisions. The reason to create a happy person is in general outweighed by the costs of creating and raising a child to the person who would do it, and so yields no obligation. The reason not to create a miserable person is not outweighed by comparable costs, since refraining from procreation is rarely costly in the way that undertaking it is. The asymmetry in obligation follows from a symmetry in value together with an asymmetry in cost.

\begin{proposition}\label{p:asym}
Given the impersonal axiology and moderate deontic principles, there is in general an obligation not to create a person whose life would not be worth living, and in general no obligation to create a person whose life would be worth living, though there is a reason to.
\end{proposition}

\begin{proof}
By Assumption~2.5 of ETH-09-a and its converse on the total view, adding a life of positive value makes an outcome better and adding a life of negative value makes it worse. The first yields a reason to create a happy person and the second a reason not to create a miserable one. A reason yields an obligation when it is not outweighed by the cost to the agent of acting on it (ETH-08-a, Section~4). The reason not to create a miserable person is acted on by refraining from procreation, whose cost is in general moderate; the reason to create a happy person is acted on by procreating, whose cost is in general great. The former therefore yields an obligation in general and the latter does not.
\end{proof}

\paragraph{The non-identity problem.} The impersonal axiology compares populations of different people, so the identity-affecting character of a choice makes no difference to its value. ETH-08 argues that obligations to future generations are owed to no one in particular and are grounded in this value; nothing in the non-identity problem then conflicts with the judgement that depleting choices are wrong.

\section{Objection and reply}

\begin{objection}
The total view implies what has been called the very repugnant conclusion: for any population of very good lives, a population containing many lives of severe suffering, together with a vast number of lives barely worth living, is better. This is a paradox the view cannot resolve by appeal to the description of lives barely worth living.
\end{objection}

\begin{reply}
The objection is right that the debunking of ETH-09-a does not apply to the suffering lives, which are what they are pictured to be. But the very repugnant conclusion combines two features each of which the total view accepts for independent reasons: that a vast number of lives barely worth living can outweigh a smaller number of very good lives, which is the repugnant conclusion, and that the suffering of some can be outweighed by sufficiently great goods in the lives of others, which every non-lexical axiology accepts in same-number cases. The case is hard because it joins two hard trade-offs, and intuition is unreliable about each of them for the reasons of scope insensitivity that ETH-09-a gives. It is a surprising consequence, not an inconsistency among principles we have reason to accept.
\end{reply}

\section{Conclusion}

Once the repugnant conclusion is accepted, the main puzzles of population ethics resolve: the mere addition paradox by rejecting its third judgement, the impossibility theorems by rejecting the condition that excludes the conclusion, the sadistic conclusion by not arising, the asymmetry by separating reasons from obligations, and the non-identity problem by the impersonality of value. The resulting field has hard cases but no paradoxes.

\begin{thebibliography}{9}\small
\bibitem[Arrhenius(2000)]{Arrhenius2000} Arrhenius, G. (2000). An impossibility theorem for welfarist axiologies. \emph{Economics and Philosophy}, 16(2), 247--266.
\bibitem[Benatar(2006)]{Benatar2006} Benatar, D. (2006). \emph{Better Never to Have Been: The Harm of Coming into Existence}. Oxford University Press.
\bibitem[Narveson(1967)]{Narveson1967} Narveson, J. (1967). Utilitarianism and new generations. \emph{Mind}, 76(301), 62--72.
\bibitem[Parfit(1984)]{Parfit1984} Parfit, D. (1984). \emph{Reasons and Persons}. Clarendon Press.
\bibitem[Zuber et~al.(2021)]{Zuber2021} Zuber, S., Venkatesh, N., T{\"a}nnsj{\"o}, T., et al. (2021). What should we agree on about the repugnant conclusion? \emph{Utilitas}, 33(4), 379--383.
\end{thebibliography}

\end{document}
